


%%%%%%%%%%%%%%%%%%%%%%%%%%%%%%%%%%%%%%%%%%%%%%%%%%%%%%%%%%%%%%
%%%%%%%%%%%%% MA-XXXX Introducción a las Demostraciones %%%%%%%%%%%%%
%%%%%%%%%%%%%%%%%%%%%%%%%%%%%%%%%%%%%%%%%%%%%%%%%%%%%%%%%%%%%%
\newpage
\subsection*{MA-0252 Introducción a las Demostraciones}
\addcontentsline{toc}{subsection}{MA-0252 Introducción a las Demostraciones}

\begin{refsection}
\begin{table}[H]
\centering{
\begin{tabular}{|ll|}
\hline
%\multicolumn{2}{|c|}{\textbf{Nivel de virtualidad del curso:} Virtual}\\ \hline
\textbf{Año:} I  & \textbf{Requisitos:} MA-0152 Matemática Exploratoria \\ 
\textbf{Ciclo:} II  & \textbf{Correquisitos:} Ninguno \\ 
\textbf{Tipo de curso:} Teórico & \textbf{Horas presenciales por semana:}   \\ 
\textbf{Créditos:} 4 & \textbf{Horas de trabajo independiente por semana:}  \\ \hline
\end{tabular}
}
\end{table}


\subsubsection*{Descripción del curso}

Este curso se encuentra en el II-ciclo de la carrera de Bach. y Lic. en Matemática y está dirigido al estudiantado que ha adquirido las habilidades y destrezas desarrolladas en el curso MA-0152 Matemática Exploratoria.

Uno de los aspectos fundamentales de este curso es que continua la transición entre la matemática a un nivel puramente operacional y la matemática más formal, la cual fue iniciada en el curso MA-0152 Matemática Exploratoria. En particular, se da énfasis al uso correcto de las bases de la Lógica Matemática para aplicarlas en la redacción de demostraciones. 

Adicionalmente a esto, durante el curso se presentan algunos fundamentos de conjuntos, relaciones y funciones, teoría de números elemental e inducción matemática, que resultan básicos en la formación matemática del estudiantado.

Durante el curso se utilizarán las intuiciones y nociones de los métodos de razonamiento inductivo y deductivo que se trabajaron en el curso de Matemática Exploratoria para desarrollar las habilidades de demostración en matem\'aticas. 




\subsubsection*{Objetivos}

\textbf{Objetivo general:} Desarrollar en el estudiantado las habilidades de comprensión y aplicación de estrategias y métodos de demostración en matemáticas.

\textbf{Objetivos específicos.} Durante el curso se espera que el estudiantado sea capaz de:

\begin{enumerate} \setlength\itemsep{1pt}
\item Comprender y apreciar la importancia de las demostraciones en matemáticas como herramientas fundamentales para establecer la veracidad de afirmaciones matemáticas.

\item Comunicar ideas matemáticas de manera clara y efectiva en forma oral y escrita.
  
\item Analizar y desglosar proposiciones matemáticas en sus componentes básicos e identificar las hipótesis o premisas.
  
\item Aplicar las técnicas básicas de lógica matemática para construir argumentos lógicos consistentes.
  
\item Aplicar las estrategias de demostración básicas, como pruebas directas, por contradicción, por contrapositiva, pruebas por casos, por inducción matemática, entre otras.
  
\item Reconocer y corregir errores en demostraciones matemáticas.

\item Aplicar los conceptos básicos sobre conjuntos, relaciones binarias y funciones y teoría de números elemental en la demostración de propiedades y teoremas matemáticos.
\end{enumerate}

\newpage
\subsubsection*{Contenidos}


\begin{enumerate}

\item \textbf{Lógica básica y estrategias de demostración.}

\begin{enumerate}
    \item Lenguaje matemático, proposiciones y formas proposicionales, conectivas lógicas (conjunción, disyunción, negación), equivalencia lógica, tautologías y contradicciones, cuantificadores, condicionales y bicondicionales. Reglas de inferencia básicas sin y con cuantificadores.
    \item Estrategias básicas de demostración: pruebas directas, por contrapositiva, por contradicción, contraejemplos, pruebas por casos, pruebas de proposiciones con cuantificadores.
\end{enumerate}

\item \textbf{Teoría de conjuntos.}

\begin{enumerate}
    \item Breve motivación sobre la importancia de la Teoría de Conjuntos elemental y su rol en la matemática moderna.
    \item Relación de pertenencia, inclusión, igualdad. 
    \item Operaciones básicas con conjuntos: unión, intersección,
    diferencia y complemento, diferencia simétrica. Propiedades básicas de estas operaciones.
    \item El conjunto potencia y el producto cartesiano.
    \item Familias de conjuntos.
\end{enumerate}


\item \textbf{Relaciones y funciones.}

\begin{enumerate}
    \item Definición y ejemplos de relaciones binarias, composición, relación inversa. 
    \item Propiedades básicas: reflexividad, simetría, transitividad y antisimetría.
    \item Definiciones y ejemplos de Orden parcial y total. 
    \item Relaciones de equivalencia, clases de equivalencia. La relación de congruencia módulo $n$ y los enteros modulo $n$.
    \item El conjunto cociente y particiones. La dualidad entre una partición y una relación de equivalencia.
    \item Definición y ejemplos de función.
    \item Imagen directa e inversa, operaciones con funciones, composición de funciones.
    \item Tipos particulares de funciones: función identidad, función característica, restricción y extensión de funciones, proyecciones de un producto cartesiano.
    \item Inversas laterales y función inversa. 
    \item Inyectividad, sobreyectividad, biyectividad y el criterio para la invertibilidad de una función.
    \item Cardinalidad y conjuntos contables.
\end{enumerate}

 

\item \textbf{Inducción matemática y recurrencia.}

\begin{enumerate}
    \item El Principio de Inducción Matemática y sus variantes.
    \item El Principio del Buen Orden.
    \item Definiciones por
    recurrencia y el Teorema de Recurrencia.
    \item Sumatorias y productorias, el factorial.
\end{enumerate}



\item \textbf{Teoría de números.}

\begin{enumerate}
    \item Divisibilidad en los números enteros, el Algoritmo de la División, 
    \item Números primos. Existencia de una infinidad de números primos y el Teorema de Factorización Única en los enteros.
    \item Máximo común divisor y mínimo común múltiplo. Propiedades básicas de estos.
    \item Ecuaciones diofánticas lineales $ax + by = c$.
    \item Congruencias.
\end{enumerate}



\end{enumerate}%Fin del enumerate de Contenidos

\subsubsection*{Metodología}

%\noindent \textbf{Lineamientos metodológicos:} La persona docente a cargo de este curso debe respetar las siguientes pautas:
%\begin{enumerate}
 %   \item 
%\end{enumerate}


\textbf{Sugerencias metodológicas:} Adicionalmente, se tienen las siguientes recomendaciones para la persona docente a cargo de este curso:
\begin{enumerate}
    \item Comienzar con ejemplos sencillos y gradualmente ir aumentando la complejidad de las demostraciones a medida que avanza el curso. Esto ayudará al estudiantado a construir una base sólida antes de enfrentarse problemas más difíciles.
    \item Fomentar la participación activa del estudiantado durante las clases. Por ejemplo, plantear preguntas, discusiones y alentar al estudiantado a presentar sus propias demostraciones en clase.
    \item Proporcionar retroalimentación regular sobre las demostraciones del estudiantado. Esto puede ser a través de correcciones escritas o discusiones en clase.
    \item Utilizar ejemplos concretos para ilustrar conceptos abstractos. Por ejemplo, al introducir un tema como la imagen y la imagen inversa de una función se puede considerar ejemplos de funciones que asignan precios a los artículos en un supermercado.
    \item Proporcionar una amplia variedad de ejercicios y problemas para que el estudiantado practique y desarrolle las habilidades para la creación y escritura correcta de demostraciones.
    \item Fomentar la colaboración entre estudiantes para resolver problemas y desarrollar demostraciones. Trabajar en grupos pequeños puede ayudar al estudiantado a discutir ideas y reforzar su comprensión.
    \item En caso de que algunos temas lo permitan, presentar al estudiantado casos de demostraciones famosas en la historia de las matemáticas para ilustrar diferentes métodos de demostración y mostrar cómo la demostración matemática ha evolucionado con el tiempo. Un ejemplo concreto de esto es la demostración de Euclides de la existencia de infinitos números primos.
\end{enumerate}

\nocite{And94}
\nocite{Blo11}
\nocite{CD}
\nocite{Dud78}
\nocite{UG11}
\nocite{Lak16}
\nocite{Mur18}
\nocite{Nic19}
\nocite{Piz02}
\nocite{Sol13}
\nocite{Ste09}
\nocite{Vel19}
\nocite{Zal14}
\nocite{Zal18}
%\subsubsection*{Bibliografía}
%El libro de Block \cite{Blo11}

\printbibliography[heading=subbibliography]
%\begin{bibunit}
%\putbib[Biblio]
%\end{bibunit}

\end{refsection}